\documentclass[10pt,a4paper]{article}

\usepackage{marktwo}
\usepackage{graphicx}

\begin{document}

\makemathtitle
    {Function-al vectors}
    {Some fourier analysis and linear algebra}
    {Vivek Kalita}
    {Department of Physics, IISER Mohali}
    {17 October 2026}

\tableofcontents

\section{Some basic Linear Algebra}
Before going into the main topic, lets first start with a basic review. 


Vector spaces (or more expressively linear spaces) are defined by two operations: addition and scalar multiplication. 

\begin{definition}[Vector Space]
Let $\mathbb{F}$ be a field. A \textbf{vector space over $\mathbb{F}$}
is an abelian group $(V,+)$ equipped with a scalar multiplication
\[
\mathbb{F}\times V \longrightarrow V,
\qquad (a,v)\longmapsto av,
\]
such that, for all $a,b\in\mathbb{F}$ and $u,v\in V$,
\begin{enumerate}
    \item \(a(u+v) = au+av,\)
    \item \((a+b)v = av+bv,\)
    \item \((ab)v = a(bv),\)
    \item \(1_{\mathbb{F}}v = v.\)
\end{enumerate}
\end{definition}

Now consider two vectors \(\mathbf u\) and \(\mathbf v\), it is easy to see that the space given by:
\[
V = \{a \mathbf{u} + b \mathbf{v} \,|\, a,b \in \mathbb F\}
\]
could generate a subspace of the vector space, such that it is a vector space in itself. 

In general, for any vector space there exist a set of vectors (the basis) whose linear combinations would produce the entire vector space.

As for example, 
\begin{enumerate}
    \item In \(\mathbb R ^3\), you can produce all the vectors with the simple basis \((1,0,0), (0,1,0)\) and \((0,0,1)\).
\end{enumerate}

\begin{definition}[Basis]
Let \(V\) be a vector space over a field \(\mathbb{F}\). A subset
\(B=\{\mathbf{e}_1,\mathbf{e}_2,\ldots,\mathbf{e}_n\}\subseteq V\)
is called a \emph{basis} of \(V\) if:

\begin{enumerate}
    \item \(B\) spans \(V\), i.e., every \(\mathbf{v}\in V\) can be
    expressed as
    \[
        \mathbf{v}
        =a_1\mathbf{e}_1+a_2\mathbf{e}_2
        +\cdots+a_n\mathbf{e}_n,
        \qquad a_i\in\mathbb{F}.
    \]

    \item \(B\) is linearly independent, i.e.,
    \[
        a_1\mathbf{e}_1+a_2\mathbf{e}_2
        +\cdots+a_n\mathbf{e}_n=\mathbf{0}
    \]
    implies
    \[
        a_1=a_2=\cdots=a_n=0.
    \]
\end{enumerate}
\end{definition}
Now till now what we've done is algebra. To induce geometry into the theory of linear spaces, we need to induce a metric.

We do this by introducing the inner product, which consequenly allows us to define the metric and also the angular measure between two vectors. 
\begin{definition}[Inner Product]
Let \(V\) be a vector space over \(\mathbb{F}\), where
\(\mathbb{F}=\mathbb{R}\) or \(\mathbb{C}\). An inner product
on \(V\) is a map
\[
    \langle\cdot,\cdot\rangle:
    V\times V\longrightarrow\mathbb{F}
\]
satisfying the following properties for all
\(\mathbf{u},\mathbf{v},\mathbf{w}\in V\) and
\(a,b\in\mathbb{F}\):

\begin{enumerate}
    \item Linearity
    \[
        \langle a\mathbf{u}+b\mathbf{v},\mathbf{w}\rangle
        =
        a\langle\mathbf{u},\mathbf{w}\rangle
        +b\langle\mathbf{v},\mathbf{w}\rangle.
    \]

    \item Conjugate symmetry
    \[
        \langle\mathbf{u},\mathbf{v}\rangle
        =
        \overline{\langle\mathbf{v},\mathbf{u}\rangle}.
    \]

    \item Positive definiteness
    \[
        \langle\mathbf{v},\mathbf{v}\rangle\geq 0,
    \]
    with equality if and only if \(\mathbf{v}=\mathbf{0}\).
\end{enumerate}
\end{definition}
The \textbf{\textcolor{blue}{norm}} of a vector \(\mathbf v\) consequently is given as: 
\[\lVert \mathbf{v} \rVert^2 = \langle \mathbf{v},\mathbf{v} \rangle.
\]

Consequently, we can also define the angle between two
non-zero vectors by
\[
    \cos\theta
    =
    \frac{\langle \mathbf{u},\mathbf{v}\rangle}
    {\lVert\mathbf{u}\rVert\,\lVert\mathbf{v}\rVert}.
\]
We could thus define a notion of a orthogonal basis. 
\begin{definition}[Orthogonal Basis]
Let \(V\) be an inner product space over \(\mathbb{F}\). A basis
\(B=\{\mathbf{e}_1,\mathbf{e}_2,\ldots,\mathbf{e}_n\}\) of
\(V\) is called an \emph{orthogonal basis} if its distinct
elements are mutually orthogonal, i.e.,
\[
    \langle\mathbf{e}_i,\mathbf{e}_j\rangle=0,
    \qquad i\neq j.
\]
If, in addition, every basis vector has unit norm,
\[
    \lVert\mathbf{e}_i\rVert=1,
    \qquad i=1,2,\ldots,n,
\]
the basis is called an \emph{orthonormal basis}.
\end{definition}

\section{Exploring functions as bases}
Consider the follwing integrals:
\[
\int_{-\pi}^{\pi}\cos(nx)\cos(mx)\,dx=\pi\delta_{nm},
\qquad
\int_{-\pi}^{\pi}\sin(nx)\sin(mx)\,dx=\pi\delta_{nm},
\qquad
\int_{-\pi}^{\pi}\sin(nx)\cos(mx)\,dx=0.
\]
Quite easily you can see where I am trying to get at. 
The functions
\[
\mathcal{B} = \{1,\, \frac{1}{\pi}\cos(nx),\, \frac{1}{\pi}\sin(nx) \}
\]
could potentially be considered as the basis of a a vector space.

\subsection{The trignometric basis}
With the basis laid over, we can define a vector space with \(\mathcal{B}\) as the basis:

\begin{definition}[Space of Trigonometric Polynomials]
Let \(\mathcal{T}\) be a vector space over \(\mathbb{R}\), spanned by
the set of functions
\[
    \mathcal{B}
    =
    \left\{
    1,\frac{1}{\pi}\cos(nx),\frac{1}{\pi}\sin(nx)
    \;\middle|\;
    n\in\mathbb{N}
    \right\}.
\]

Thus, every \(f\in\mathcal{T}\) can be expressed as
\[
    f(x)
    =
    a_0+
    \sum_{n=1}^{N}
    \left[
    \frac{a_n}{\pi}\cos(nx)
    +
    \frac{b_n}{\pi}\sin(nx)
    \right],
\]
where \(a_n,b_n\in\mathbb{R}\) and \(N\) is finite. We equip \(\mathcal{T}\) with the inner product
\[
    \langle f,g\rangle
    =
    \int_{-\pi}^{\pi}f(x)g(x)\,dx,
    \qquad f,g\in\mathcal{T}.
\]
\end{definition}
In general, given the function \(f(x)\), we can also extract any of the coefficients. And also we have some very basic results:

\begin{theorem}[Basic Results]
Let \(f\in\mathcal{T}\) have the representation
\[
    f(x)=a_0+\sum_{n=1}^{N}
    \left[
        \frac{a_n}{\pi}\cos(nx)
        +\frac{b_n}{\pi}\sin(nx)
    \right].
\]
Then the following results hold:

\begin{enumerate}
    \item \textup{Fourier Coefficients:}  The coefficients are given by
    \[
        a_0=\frac{1}{2\pi}\int_{-\pi}^{\pi}f(x)\,dx,
        \quad
        a_n=\int_{-\pi}^{\pi}f(x)\cos(nx)\,dx,
        \quad
        b_n=\int_{-\pi}^{\pi}f(x)\sin(nx)\,dx.
    \]

    \item \textup{Parseval's Identity:} The squared norm of \(f\) is given by
    \[
        \lVert f\rVert^2
        =\langle f,f\rangle
        =\int_{-\pi}^{\pi}f^2(x)\,dx
        =2\pi a_0^2+\frac{1}{\pi}
        \sum_{n=1}^{N}(a_n^2+b_n^2).
    \]
\end{enumerate}
\end{theorem}

\subsection{When can you fourier-expand?}
With our choice of basis \(\mathcal{B}\), the Fourier coefficients
are given by
\[
    a_0=\frac{1}{2\pi}\int_{-\pi}^{\pi}f(x)\,dx,
    \qquad
    a_n=\int_{-\pi}^{\pi}f(x)\cos(nx)\,dx,
    \qquad
    b_n=\int_{-\pi}^{\pi}f(x)\sin(nx)\,dx.
\]

For all these integrals to exist, it is sufficient that
\(f\in L^1([-\pi,\pi])\), which means that \(f(x)\) is
absolutely integrable over the interval, i.e.,
\[
    \int_{-\pi}^{\pi}|f(x)|\,dx<\infty.
\]

To see why this condition is sufficient, observe that
\[
    |\cos(nx)|\leq 1,
    \qquad
    |\sin(nx)|\leq 1.
\]
Consequently,
\[
\begin{aligned}
    \left|\int_{-\pi}^{\pi}f(x)\cos(nx)\,dx\right|
    &\leq \int_{-\pi}^{\pi}|f(x)\cos(nx)|\,dx\\
    &\leq \int_{-\pi}^{\pi}|f(x)|\,dx<\infty.
\end{aligned}
\]
The same argument applies to the sine coefficients. Thus,
all the Fourier coefficients are well-defined.

\begin{theorem}[Existence of Fourier Coefficients]
The Fourier coefficients of a function \(f\) are given by
\[
    A_0=\frac{1}{2\pi}\int_{-\pi}^{\pi}f(x)\,dx,
    \qquad
    A_n=\int_{-\pi}^{\pi}f(x)\cos(nx)\,dx,
    \qquad
    B_n=\int_{-\pi}^{\pi}f(x)\sin(nx)\,dx,
\]
provided that \(f\) is absolutely integrable over the
interval \([-\pi,\pi]\). That is: 
\[
    \int_{-\pi}^{\pi}|f(x)|\,dx<\infty.
\]
\end{theorem}
\begin{proof}
\begin{enumerate}
    \item (\(\Rightarrow\)) Suppose all the Fourier coefficients exist as
finite integrals. In particular, \(A_0\) exists, so
\[
    \int_{-\pi}^{\pi}f(x)\,dx
\]
is finite. Hence \(f\) is absolutely integrable, i.e.,
\[
    \int_{-\pi}^{\pi}|f(x)|\,dx<\infty.
\]
\item (\(\Leftarrow\)) Conversely, suppose that
\[
    \int_{-\pi}^{\pi}|f(x)|\,dx<\infty.
\]
Since
\[
    |\cos(nx)|\leq 1,
    \qquad
    |\sin(nx)|\leq 1,
\]
we have
\[
\begin{aligned}
    \int_{-\pi}^{\pi}|f(x)\cos(nx)|\,dx
    &\leq \int_{-\pi}^{\pi}|f(x)|\,dx<\infty,\\
    \int_{-\pi}^{\pi}|f(x)\sin(nx)|\,dx
    &\leq \int_{-\pi}^{\pi}|f(x)|\,dx<\infty.
\end{aligned}
\]
Thus, all the Fourier coefficients exist as finite 
integrals.
\end{enumerate}
\end{proof}

\begin{example}
But now, consider the function
\[
    f(x)=\frac{1}{\sqrt{|x|}},
    \qquad x\in[-\pi,\pi],\quad x\neq 0.
\]
\begin{enumerate}
    \item If we define \(f(0)=0\), then we can easily verify that
\[
    \int_{-\pi}^{\pi}|f(x)|\,dx
    =4\sqrt{\pi}<\infty.
\]
Thus, \(f\in L^1([-\pi,\pi])\).
\item However, it is easily seen that
\[
    |f(x)|^2=\frac{1}{|x|},
\]
and hence
\[
    \int_{-\pi}^{\pi}|f(x)|^2\,dx
    =2\int_0^\pi\frac{1}{x}\,dx
    =\infty.
\]
Therefore, \(f\notin L^2([-\pi,\pi])\).
\end{enumerate}
\end{example}


Essentially, this demonstrates that absolute integrability alone is
not sufficient to guarantee a finite norm. Consequently, if
we wish to equip our function space with the inner product
\[
    \langle f,g\rangle
    =\int_{-\pi}^{\pi}f(x)g(x)\,dx,
\]
we must restrict ourselves to square-integrable functions.
Thus, our function space should be \(L^2([-\pi,\pi])\),
which ensures that every function has a finite norm and that
its Fourier coefficients are well-defined.

With this final restriction, we now have a complete description of the trignometric space.
\begin{theorem}[Fourier Expansion in \(L^2\)]
The trigonometric system
\[
    \mathcal{B}
    =
    \left\{
    1,\frac{1}{\pi}\cos(nx),\frac{1}{\pi}\sin(nx)
    \;\middle|\; n\in\mathbb{N}
    \right\}
\]
forms a complete orthogonal basis of
\(L^2([-\pi,\pi])\). Consequently, every \(f\in L^2([-\pi,\pi])\) admits a
Fourier expansion
\[
    f(x)=a_0+\sum_{n=1}^{\infty}
    \left[
        \frac{a_n}{\pi}\cos(nx)
        +\frac{b_n}{\pi}\sin(nx)
    \right],
\]
where the coefficients are given by the corresponding
Fourier integrals, and the partial sums converge to \(f\)
in the \(L^2\) norm.
\end{theorem}
The formal proof of this would require piecing a few more parts. And hence would be discussed later.


There is yet another very interesting thing to observe about the fourier coefficients. The zeros of \(\cos(nx)\) are located at
\[
    x_k=\frac{(2k+1)\pi}{2n},
    \qquad k\in\mathbb Z.
\]
The separation between two consecutive zeros is
\[
    \Delta x=\frac{\pi}{n}\longrightarrow 0
    \qquad \text{as } n\to\infty.
\]

Approximate \(f(x)\) by a step function, taking a
constant value \(c_j\) over each interval
\([p_j,q_j]\). Then
\[
    b_n=\int_{-\pi}^{\pi}f(x)\cos(nx)\,dx
    \approx
    \sum_{j=1}^{m}c_j
    \int_{p_j}^{q_j}\cos(nx)\,dx.
\]

Since
\[
    \left|\int_{p_j}^{q_j}\cos(nx)\,dx\right|
    \leq \frac{2}{n},
\]
we obtain
\[
    |b_n|
    \lesssim
    \frac{2}{n}\sum_{j=1}^{m}|c_j|
    \longrightarrow 0
    \qquad \text{as }n\to\infty.
\]

A similar argument holds for the other Fourier coefficient.
\begin{theorem}[Vanishing of Fourier Coefficients]
Let \(f\in L^1([-\pi,\pi])\) have the Fourier expansion
\[
    f(x)=a_0+\sum_{n=1}^{\infty}
    \left[
        \frac{a_n}{\pi}\cos(nx)
        +\frac{b_n}{\pi}\sin(nx)
    \right].
\]
Then the Fourier coefficients satisfy
\[
    \lim_{n\to\infty}a_n=0,
    \qquad
    \lim_{n\to\infty}b_n=0.
\]
\end{theorem}

\subsection{The complex exponential as the basis}
Consider the following integral:
\[
    \int_{-\pi}^{\pi}e^{imx}e^{-inx}\,dx
    =
    2\pi\delta_{mn}.
\]
This isn't suprising, as the complex exponential inherently contains the trignometric terms. In general you can define the vector space:
\begin{definition}[Complex Fourier Series]
Let \(\mathcal{H}=L^2([-\pi,\pi];\mathbb{C})\) be the
vector space of complex-valued square-integrable functions,
equipped with the inner product
\[
    \langle f,g\rangle
    =
    \int_{-\pi}^{\pi}f(x)g^*(x)\,dx,
    \qquad f,g\in\mathcal{H}.
\]

The set of complex exponential functions
\[
    \mathcal{B}
    =
    \left\{
    e^{imx}\mid m\in\mathbb{Z}
    \right\}
\]
forms a complete orthogonal basis of \(\mathcal{H}\). Consequently, every \(f\in\mathcal{H}\) admits a
Fourier expansion of the form
\[
    f(x)=\sum_{m=-\infty}^{\infty}c_m e^{imx},
\]
where the Fourier coefficients are given by
\[
    c_m=\frac{1}{2\pi}
    \int_{-\pi}^{\pi}f(x)e^{-imx}\,dx.
\]
\end{definition}

We can derive some basic results from the complex representation.

\begin{enumerate}
    \item \textcolor{blue}{Parseval's Identity:}
    \[
        \lVert f\rVert^2
        =2\pi\sum_{m=-\infty}^{\infty}c_m c_m^{*}.
    \]

    \item \textcolor{blue}{Riemann-Lebesgue Lemma:}

    For \(f\in L^1([-\pi,\pi])\), its Fourier coefficients satisfy
    \[
        \lim_{|m|\to\infty}c_m=0.
    \]

    \item \textcolor{blue}{Conjugate Symmetry:}

    If \(f(x)\) is real-valued, its Fourier coefficients satisfy
    \[
        c_{-m}=c_m^{*}.
    \]

    and they would be related to the sine and cosine coefficient (for the real-valued fourier expansion) as:

    \[
        a_0=c_0,\qquad
        a_n=\pi(c_n+c_{-n}),\qquad
        b_n=i\pi(c_n-c_{-n}).
    \]

    \item \textcolor{blue}{Inner Product Identity:}

    For \(f,g\in L^2([-\pi,\pi])\), with Fourier coefficients
    \(c_m\) and \(d_m\), respectively,
    \[
        \langle f,g\rangle
        =\int_{-\pi}^{\pi}f(x)g^{*}(x)\,dx
        =2\pi\sum_{m=-\infty}^{\infty}c_m d_m^{*}.
    \]
\end{enumerate}

\end{document}